\subsection{Removing the cubic height without asymptotics}\label{sec:airyformal}

The extra term $T^3$ can be dealt with algebraically. This avoids
using an asymptotic formula for a cubic Gaussian integral.

\begin{lemma}[An algebraic Riccati obstruction]\label{lem:airyriccati}
The equation $f_x=x-f^2$ has no solution algebraic over $\C(x)$.
\end{lemma}
\begin{proof}
At a finite point, a finite branch of $f$ satisfies a regular analytic
first-order equation and cannot branch. At a pole, put $g=1/f$.
Then $g_x=1-xg^2$, again a regular equation; a zero of $g$ is simple.
Thus the normalized algebraic curve of $\C(x,f)$ is unramified over
every finite $x$. If its degree is $d$ and there are $r$ points over
infinity, Riemann--Hurwitz gives $2g-2=-d-r$. This forces $d=r=1$.
So $f$ would be rational. An integer-order pole of $f$ at infinity
makes $f^2$ dominate both $x$ and $f_x$; if $f$ is bounded there,
the term $x$ cannot cancel. Both possibilities are impossible.
\end{proof}

\begin{lemma}[Cubic Gaussian obstruction]\label{lem:cubicgaussian}
Let $A,B,C\in\C(z)$ be regular at zero, with $A(0)=C(0)=0$,
$B(0)=-1/2$ and $A\ne0$. Define the formal unit
\[
 I=\E_T\exp\bigl(AT^3+(B+1/2)T^2+CT\bigr).
\]
Then $I'/I$ is not rational in $z$.
\end{lemma}
\begin{proof}
Write $M_j=\E[T^j\exp(AT^3+(B+1/2)T^2+CT)]$ and $J=M_1$.
Gaussian integration by parts gives, for $j\geq0$,
\begin{equation}\label{eq:momentrec}
 3A M_{j+2}+2B M_{j+1}+C M_j+jM_{j-1}=0,
\end{equation}
where the last term is zero for $j=0$. Also
$I'=A'M_3+B'M_2+C'J$ and $J'=A'M_4+B'M_3+C'M_2$.
For clarity, the first two eliminations are
\[
 M_2=-\frac{2BJ+CI}{3A},\qquad
 M_3=\frac{(4B^2-3AC)J+(2BC-3A)I}{9A^2}.
\]
The case $j=2$ of \eqref{eq:momentrec} supplies $M_4$.

In an algebraic extension choose $a^3=3A$ and put
\[
 X=\frac{B^2-3AC}{a^4},\qquad F=\frac JI,
 \qquad f=-aF-\frac B{a^2}.
\]
Substitution in these three recurrences, using $a'=A'/a^2$, gives
\begin{align}
 \frac{I'}I&=U-aX'F,
 &U&=\frac{A'(2BC-3A)}{9A^2}-\frac{B'C}{3A},\label{eq:airyfirst}\\
 f'&=X'(X-f^2).\label{eq:airypullback}
\end{align}
These identities follow by differentiation and polynomial substitution;
they do not involve an integration contour or a limit of a Gaussian
integral. They can also be checked by substituting $A=a^3/3$ in
\eqref{eq:momentrec} before differentiating.

If $A$ has order $q>0$, then $X$ has a pole of order $4q/3$,
because $B^2-3AC$ is a unit. Thus $X'\ne0$ and $aX'$ is a nonzero
rational function. If $I'/I$ were rational, \eqref{eq:airyfirst}
would make $F$ rational and $f$ algebraic over $\C(z)$. The function
field $\C(z,a)$ is finite over $\C(X)$, so $f$ would be algebraic
over $\C(X)$. Equation~\eqref{eq:airypullback} contradicts
Lemma~\ref{lem:airyriccati}.
\end{proof}

